% Accuracy subsection removed from paper/epi_int_3d.tex (Section 2.4) on 2026-10-17; kept here
% for possible reuse (e.g. in a later section or appendix). Restore with % Accuracy subsection removed from paper/epi_int_3d.tex (Section 2.4) on 2026-10-17; kept here
% for possible reuse (e.g. in a later section or appendix). Restore with % Accuracy subsection removed from paper/epi_int_3d.tex (Section 2.4) on 2026-10-17; kept here
% for possible reuse (e.g. in a later section or appendix). Restore with % Accuracy subsection removed from paper/epi_int_3d.tex (Section 2.4) on 2026-10-17; kept here
% for possible reuse (e.g. in a later section or appendix). Restore with \input{unused/accuracy}.
% Labels used here: sec:accuracy, eqn:residual, tab:accuracy. It refers to eqn:zvz, eqn:nu_r,
% eqn:eom_expand, eqn:Phi3d and sec:appendix_inplane in the main file.
% Numbers in Table 1 come from notes/numerics/vertical_motion_check.py; the residual is
% sympy check V6 in notes/sympy/vertical_motion.py.

\subsection{Accuracy}
\label{sec:accuracy}

Eqns.\ (\ref{eqn:zvz}) neglect the variation of $\nu^2(r)$ over the particle's radial
excursion. Inserting $r = a(1 + \delta)$ with $\delta = {\cal O}(e)$ into Eqn.\ (\ref{eqn:nu_r})
gives $\nu^2(r) \simeq \nu_a^2 + a\delta\,d\nu^2/da$, so Eqns.\ (\ref{eqn:zvz}) satisfy the exact
vertical equation of motion $\ddot{z} = -\partial\Phi/\partial z$ with a residual of order
${\cal O}(ei)$,
\begin{equation}
    \label{eqn:residual}
    \ddot{z} + \nu_a^2z \simeq \frac{3GM_p}{a^3}\left[1 + \frac{15}{2}J_2\left(\frac{R_p}{a}\right)^2\right]\delta z,
\end{equation}
plus terms of order ${\cal O}(i^3)$. Likewise the in-plane solution ignores the ${\cal O}(z^2)$
term in the radial Eqn.\ (\ref{eqn:eom_expand}), {\it i.e.}\ the ${\cal O}(i^2)$ change in the
radial force on a particle at height $z$. The solution is therefore accurate to first order,
with errors of order ${\cal O}(e^2, i^2, ei)$, the same order as the in-plane solution alone.
All equations in this section and in Appendix \ref{sec:appendix_inplane} were verified
symbolically.

Table \ref{tab:accuracy} confirms this numerically for Saturn-like $J_2 = 0.016$ and $a = 2R_p$.
Converting $10^5$ random particles having $e, i \le 0.05$ from coordinates to elements and
back reproduces the coordinates with relative errors $\lesssim 10^{-15}$, and $10^4$ repeated
round trips change $r$ by $\lesssim 2\times10^{-15}a$ and $z$ by $\lesssim 2\times10^{-13}a$.
After one orbit, the position error relative to an exact integration of the particle's motion
in potential (\ref{eqn:Phi3d}) falls as the square of $e$ and $i$.

\begin{deluxetable}{ccc}
\tablecaption{Position error after one orbit\label{tab:accuracy}}
\tablehead{\colhead{$e$} & \colhead{$i$} & \colhead{$|\Delta\mathbf{r}|/a$}}
\startdata
0.02  & 0     & $3.7\times10^{-3}$ \\
0.01  & 0     & $9.3\times10^{-4}$ \\
0     & 0.02  & $4.0\times10^{-3}$ \\
0     & 0.01  & $9.9\times10^{-4}$ \\
0.02  & 0.02  & $7.7\times10^{-3}$ \\
0.01  & 0.01  & $1.9\times10^{-3}$ \\
0.005 & 0.005 & $4.8\times10^{-4}$ \\
\enddata
\tablecomments{Error in the drift-step position of a particle with $a=2R_p$ and $J_2=0.016$
after one orbit, relative to an exact numerical integration in the potential
(\ref{eqn:Phi3d}).}
\end{deluxetable}

\cite{BL94} and \cite{RS06} give conversions between coordinates and geometric
elements for inclined orbits about an oblate planet that are accurate to second order in $e$
and $i$. But those conversions are not exact inverses of each other (\citealt{H25};
the inversion of \citealt{RS06} is iterative and accurate to second order), so the repeated
conversions made during a long simulation accumulate error. The first-order equations given here trade that higher accuracy
for exact reversibility, as \cite{H25} did for the in-plane motion.
.
% Labels used here: sec:accuracy, eqn:residual, tab:accuracy. It refers to eqn:zvz, eqn:nu_r,
% eqn:eom_expand, eqn:Phi3d and sec:appendix_inplane in the main file.
% Numbers in Table 1 come from notes/numerics/vertical_motion_check.py; the residual is
% sympy check V6 in notes/sympy/vertical_motion.py.

\subsection{Accuracy}
\label{sec:accuracy}

Eqns.\ (\ref{eqn:zvz}) neglect the variation of $\nu^2(r)$ over the particle's radial
excursion. Inserting $r = a(1 + \delta)$ with $\delta = {\cal O}(e)$ into Eqn.\ (\ref{eqn:nu_r})
gives $\nu^2(r) \simeq \nu_a^2 + a\delta\,d\nu^2/da$, so Eqns.\ (\ref{eqn:zvz}) satisfy the exact
vertical equation of motion $\ddot{z} = -\partial\Phi/\partial z$ with a residual of order
${\cal O}(ei)$,
\begin{equation}
    \label{eqn:residual}
    \ddot{z} + \nu_a^2z \simeq \frac{3GM_p}{a^3}\left[1 + \frac{15}{2}J_2\left(\frac{R_p}{a}\right)^2\right]\delta z,
\end{equation}
plus terms of order ${\cal O}(i^3)$. Likewise the in-plane solution ignores the ${\cal O}(z^2)$
term in the radial Eqn.\ (\ref{eqn:eom_expand}), {\it i.e.}\ the ${\cal O}(i^2)$ change in the
radial force on a particle at height $z$. The solution is therefore accurate to first order,
with errors of order ${\cal O}(e^2, i^2, ei)$, the same order as the in-plane solution alone.
All equations in this section and in Appendix \ref{sec:appendix_inplane} were verified
symbolically.

Table \ref{tab:accuracy} confirms this numerically for Saturn-like $J_2 = 0.016$ and $a = 2R_p$.
Converting $10^5$ random particles having $e, i \le 0.05$ from coordinates to elements and
back reproduces the coordinates with relative errors $\lesssim 10^{-15}$, and $10^4$ repeated
round trips change $r$ by $\lesssim 2\times10^{-15}a$ and $z$ by $\lesssim 2\times10^{-13}a$.
After one orbit, the position error relative to an exact integration of the particle's motion
in potential (\ref{eqn:Phi3d}) falls as the square of $e$ and $i$.

\begin{deluxetable}{ccc}
\tablecaption{Position error after one orbit\label{tab:accuracy}}
\tablehead{\colhead{$e$} & \colhead{$i$} & \colhead{$|\Delta\mathbf{r}|/a$}}
\startdata
0.02  & 0     & $3.7\times10^{-3}$ \\
0.01  & 0     & $9.3\times10^{-4}$ \\
0     & 0.02  & $4.0\times10^{-3}$ \\
0     & 0.01  & $9.9\times10^{-4}$ \\
0.02  & 0.02  & $7.7\times10^{-3}$ \\
0.01  & 0.01  & $1.9\times10^{-3}$ \\
0.005 & 0.005 & $4.8\times10^{-4}$ \\
\enddata
\tablecomments{Error in the drift-step position of a particle with $a=2R_p$ and $J_2=0.016$
after one orbit, relative to an exact numerical integration in the potential
(\ref{eqn:Phi3d}).}
\end{deluxetable}

\cite{BL94} and \cite{RS06} give conversions between coordinates and geometric
elements for inclined orbits about an oblate planet that are accurate to second order in $e$
and $i$. But those conversions are not exact inverses of each other (\citealt{H25};
the inversion of \citealt{RS06} is iterative and accurate to second order), so the repeated
conversions made during a long simulation accumulate error. The first-order equations given here trade that higher accuracy
for exact reversibility, as \cite{H25} did for the in-plane motion.
.
% Labels used here: sec:accuracy, eqn:residual, tab:accuracy. It refers to eqn:zvz, eqn:nu_r,
% eqn:eom_expand, eqn:Phi3d and sec:appendix_inplane in the main file.
% Numbers in Table 1 come from notes/numerics/vertical_motion_check.py; the residual is
% sympy check V6 in notes/sympy/vertical_motion.py.

\subsection{Accuracy}
\label{sec:accuracy}

Eqns.\ (\ref{eqn:zvz}) neglect the variation of $\nu^2(r)$ over the particle's radial
excursion. Inserting $r = a(1 + \delta)$ with $\delta = {\cal O}(e)$ into Eqn.\ (\ref{eqn:nu_r})
gives $\nu^2(r) \simeq \nu_a^2 + a\delta\,d\nu^2/da$, so Eqns.\ (\ref{eqn:zvz}) satisfy the exact
vertical equation of motion $\ddot{z} = -\partial\Phi/\partial z$ with a residual of order
${\cal O}(ei)$,
\begin{equation}
    \label{eqn:residual}
    \ddot{z} + \nu_a^2z \simeq \frac{3GM_p}{a^3}\left[1 + \frac{15}{2}J_2\left(\frac{R_p}{a}\right)^2\right]\delta z,
\end{equation}
plus terms of order ${\cal O}(i^3)$. Likewise the in-plane solution ignores the ${\cal O}(z^2)$
term in the radial Eqn.\ (\ref{eqn:eom_expand}), {\it i.e.}\ the ${\cal O}(i^2)$ change in the
radial force on a particle at height $z$. The solution is therefore accurate to first order,
with errors of order ${\cal O}(e^2, i^2, ei)$, the same order as the in-plane solution alone.
All equations in this section and in Appendix \ref{sec:appendix_inplane} were verified
symbolically.

Table \ref{tab:accuracy} confirms this numerically for Saturn-like $J_2 = 0.016$ and $a = 2R_p$.
Converting $10^5$ random particles having $e, i \le 0.05$ from coordinates to elements and
back reproduces the coordinates with relative errors $\lesssim 10^{-15}$, and $10^4$ repeated
round trips change $r$ by $\lesssim 2\times10^{-15}a$ and $z$ by $\lesssim 2\times10^{-13}a$.
After one orbit, the position error relative to an exact integration of the particle's motion
in potential (\ref{eqn:Phi3d}) falls as the square of $e$ and $i$.

\begin{deluxetable}{ccc}
\tablecaption{Position error after one orbit\label{tab:accuracy}}
\tablehead{\colhead{$e$} & \colhead{$i$} & \colhead{$|\Delta\mathbf{r}|/a$}}
\startdata
0.02  & 0     & $3.7\times10^{-3}$ \\
0.01  & 0     & $9.3\times10^{-4}$ \\
0     & 0.02  & $4.0\times10^{-3}$ \\
0     & 0.01  & $9.9\times10^{-4}$ \\
0.02  & 0.02  & $7.7\times10^{-3}$ \\
0.01  & 0.01  & $1.9\times10^{-3}$ \\
0.005 & 0.005 & $4.8\times10^{-4}$ \\
\enddata
\tablecomments{Error in the drift-step position of a particle with $a=2R_p$ and $J_2=0.016$
after one orbit, relative to an exact numerical integration in the potential
(\ref{eqn:Phi3d}).}
\end{deluxetable}

\cite{BL94} and \cite{RS06} give conversions between coordinates and geometric
elements for inclined orbits about an oblate planet that are accurate to second order in $e$
and $i$. But those conversions are not exact inverses of each other (\citealt{H25};
the inversion of \citealt{RS06} is iterative and accurate to second order), so the repeated
conversions made during a long simulation accumulate error. The first-order equations given here trade that higher accuracy
for exact reversibility, as \cite{H25} did for the in-plane motion.
.
% Labels used here: sec:accuracy, eqn:residual, tab:accuracy. It refers to eqn:zvz, eqn:nu_r,
% eqn:eom_expand, eqn:Phi3d and sec:appendix_inplane in the main file.
% Numbers in Table 1 come from notes/numerics/vertical_motion_check.py; the residual is
% sympy check V6 in notes/sympy/vertical_motion.py.

\subsection{Accuracy}
\label{sec:accuracy}

Eqns.\ (\ref{eqn:zvz}) neglect the variation of $\nu^2(r)$ over the particle's radial
excursion. Inserting $r = a(1 + \delta)$ with $\delta = {\cal O}(e)$ into Eqn.\ (\ref{eqn:nu_r})
gives $\nu^2(r) \simeq \nu_a^2 + a\delta\,d\nu^2/da$, so Eqns.\ (\ref{eqn:zvz}) satisfy the exact
vertical equation of motion $\ddot{z} = -\partial\Phi/\partial z$ with a residual of order
${\cal O}(ei)$,
\begin{equation}
    \label{eqn:residual}
    \ddot{z} + \nu_a^2z \simeq \frac{3GM_p}{a^3}\left[1 + \frac{15}{2}J_2\left(\frac{R_p}{a}\right)^2\right]\delta z,
\end{equation}
plus terms of order ${\cal O}(i^3)$. Likewise the in-plane solution ignores the ${\cal O}(z^2)$
term in the radial Eqn.\ (\ref{eqn:eom_expand}), {\it i.e.}\ the ${\cal O}(i^2)$ change in the
radial force on a particle at height $z$. The solution is therefore accurate to first order,
with errors of order ${\cal O}(e^2, i^2, ei)$, the same order as the in-plane solution alone.
All equations in this section and in Appendix \ref{sec:appendix_inplane} were verified
symbolically.

Table \ref{tab:accuracy} confirms this numerically for Saturn-like $J_2 = 0.016$ and $a = 2R_p$.
Converting $10^5$ random particles having $e, i \le 0.05$ from coordinates to elements and
back reproduces the coordinates with relative errors $\lesssim 10^{-15}$, and $10^4$ repeated
round trips change $r$ by $\lesssim 2\times10^{-15}a$ and $z$ by $\lesssim 2\times10^{-13}a$.
After one orbit, the position error relative to an exact integration of the particle's motion
in potential (\ref{eqn:Phi3d}) falls as the square of $e$ and $i$.

\begin{deluxetable}{ccc}
\tablecaption{Position error after one orbit\label{tab:accuracy}}
\tablehead{\colhead{$e$} & \colhead{$i$} & \colhead{$|\Delta\mathbf{r}|/a$}}
\startdata
0.02  & 0     & $3.7\times10^{-3}$ \\
0.01  & 0     & $9.3\times10^{-4}$ \\
0     & 0.02  & $4.0\times10^{-3}$ \\
0     & 0.01  & $9.9\times10^{-4}$ \\
0.02  & 0.02  & $7.7\times10^{-3}$ \\
0.01  & 0.01  & $1.9\times10^{-3}$ \\
0.005 & 0.005 & $4.8\times10^{-4}$ \\
\enddata
\tablecomments{Error in the drift-step position of a particle with $a=2R_p$ and $J_2=0.016$
after one orbit, relative to an exact numerical integration in the potential
(\ref{eqn:Phi3d}).}
\end{deluxetable}

\cite{BL94} and \cite{RS06} give conversions between coordinates and geometric
elements for inclined orbits about an oblate planet that are accurate to second order in $e$
and $i$. But those conversions are not exact inverses of each other (\citealt{H25};
the inversion of \citealt{RS06} is iterative and accurate to second order), so the repeated
conversions made during a long simulation accumulate error. The first-order equations given here trade that higher accuracy
for exact reversibility, as \cite{H25} did for the in-plane motion.
